% Inference/Model

% 1. Use duckdb to read relbench-style dataset
% 2. For a specific inference (x, t) pair construct  a sequence of s = emb(x, t_i), y_i pair  of length k from the duckdb database. Here if i<j, then t_i < t_j < t. emb will be a tgnn/graph_transformer/relational encoder. For now tgat and GT encoder, later we will use tgn encoder. 
% 3. Feed [s, emb(x,t)] to an LLM and predict Y


% Training Plan:
% 1. Keep the llm frozen. No fine-tuning of LLM
% 2. Pick a set of relbench style datasets
% 3. Pick all the tasks that can be defined on it.
% 4. Randomly sample a batch of train (x, t), y pairs from all the tasks and all the datasets. Sampler should sample without replacement in and across batch so, for one train epoch each train x is just seen once. 
% 5. Objective to mix different dataset and task in the same batch is to train tgnn encoder in a task agnostic way. so, the embedding is task agnostic. Before passing to llm, as each sequence is for one specific task, we will add a task description, then the sequence of train icl examples, and target emb(x, t). So, a batch will be a collection of different tasks but a sequence will have example for just one task and one entity.
% 6. Loss and backprop will only update encoder params, treating LLM as a fixed function.




\documentclass[11pt]{article}

\usepackage[T1]{fontenc}
\usepackage{textcomp}
\usepackage{mathptmx}
\usepackage{amsmath, amssymb}
\usepackage{graphicx}
\usepackage{booktabs}
\usepackage{enumitem}
\usepackage{geometry}
\geometry{margin=1in}

\title{\textbf{Relational Foundation Models via Temporal In-Context Learning over Graph Encoders}}
\author{}
\date{}

\begin{document}

\maketitle

\section{Introduction}

We propose a novel framework for \textit{Relational Foundation Models (RFM)} that combines temporal relational encoders with in-context learning (ICL) using a frozen Large Language Model (LLM) \cite{dummy-rfm-note}. 

The core idea is to learn a \textbf{task-agnostic relational encoder} whose outputs can be interpreted by a frozen LLM through example-based conditioning. This design decouples representation learning from task-specific reasoning, enabling modularity, generalization, and adaptability to future LLMs.

\section{Problem Setup}

We assume a relational dataset stored in a database (e.g., DuckDB) with temporal information. Each query is of the form:

\begin{equation}
(x, t) \rightarrow y
\end{equation}

where $x$ is an entity, $t$ is a timestamp, and $y$ is the target variable. Tasks include link prediction, node classification/regression, and forecasting.

\section{Framework Overview}

\subsection{Context Construction}

For each query $(x,t)$, we construct a temporally valid context:

\begin{equation}
\mathcal{C} = \{(x, t_i, y_i)\}_{i=1}^{k}, \quad t_i < t
\end{equation}

\subsection{Encoding}

A temporal relational encoder produces embeddings:

\begin{equation}
h_i = f_\theta(x, t_i), \quad h_q = f_\theta(x, t)
\end{equation}

\subsection{Projection and Interface}

To align with LLM input space, we introduce a projection layer:

\begin{equation}
z_i = W h_i, \quad z_q = W h_q
\end{equation}

\subsection{ICL Sequence Construction}

The LLM receives:

\begin{equation}
[(z_i, y_i)]_{i=1}^{k}, \quad z_q
\end{equation}

and predicts:

\begin{equation}
\hat{y} = \text{LLM}(\text{Prompt}, \text{Sequence})
\end{equation}

\section{Training Strategy}

\subsection{Frozen LLM}

The LLM is kept fixed. Only the encoder $f_\theta$ and projection $W$ are trained.

\subsection{Multi-Task, Multi-Dataset Training}

Training batches mix:
\begin{itemize}
\item multiple datasets
\item multiple tasks
\item diverse $(x,t)$ samples without repetition per epoch
\end{itemize}

This enforces task-agnostic representation learning.

\subsection{Objective Function}

Primary loss:
\begin{equation}
\mathcal{L}_{LLM} = \ell(\hat{y}, y)
\end{equation}

Auxiliary loss (optional but recommended):
\begin{equation}
\mathcal{L}_{aux} = \ell(g(h_q), y)
\end{equation}

Final objective:
\begin{equation}
\mathcal{L} = \mathcal{L}_{LLM} + \lambda \mathcal{L}_{aux}
\end{equation}

\section{Embedding--LLM Interface as an Alignment Problem}

A central challenge is aligning encoder outputs with LLM input expectations.

\subsection{Types of Misalignment}

\begin{itemize}
\item \textbf{Dimensional Alignment:} $h \in \mathbb{R}^d$ vs token embeddings in $\mathbb{R}^{d_{LLM}}$
\item \textbf{Distributional Alignment:} encoder outputs do not follow token embedding distribution
\item \textbf{Semantic Alignment:} embeddings lack inherent meaning in LLM space
\end{itemize}

\subsection{Implications}

Even with projection $W$, the LLM receives vectors that do not correspond to meaningful tokens. Therefore, the encoder must learn representations that are \textbf{interpretable under the fixed LLM}.

\subsection{Design Principle}

We define the goal as learning:
\begin{equation}
h = f_\theta(x,t)
\end{equation}
such that:
\begin{equation}
\text{LLM}(\text{format}(h, \mathcal{C})) \rightarrow y
\end{equation}

\section{Design Choices and Alternatives}

\subsection{Encoder Choices}

\begin{center}
\begin{tabular}{lcc}
\toprule
Encoder & Pros & Cons \\
\midrule
TGAT & Efficient, simple & Limited memory \\
Graph Transformer & Expressive & High cost \\
TGN & Temporal memory & Complex \\
\bottomrule
\end{tabular}
\end{center}

\subsection{Context Construction Strategies}

\begin{itemize}
\item Full history (expensive)
\item Recent-$k$ history (efficient)
\item Graph-aware sampling (informative but costly)
\item Hybrid (recommended)
\end{itemize}

\subsection{Embedding Interface Strategies}

\begin{itemize}
\item Text serialization (simple, weak)
\item Linear projection (recommended baseline)
\item Discretization (token-like but lossy)
\item Cross-attention adapter (powerful but intrusive)
\end{itemize}

\section{Advantages}

\begin{itemize}
\item \textbf{Task Generalization:} single encoder across tasks
\item \textbf{Modularity:} decouples encoder and reasoning
\item \textbf{Temporal Reasoning:} naturally incorporates history
\item \textbf{Future Compatibility:} enables plug-and-play with future LLMs
\end{itemize}

\section{Limitations and Risks}

\begin{itemize}
\item \textbf{Alignment Difficulty:} embeddings may not align with LLM expectations
\item \textbf{Weak Gradient Signal:} supervision passes through frozen LLM
\item \textbf{Sampling Bottleneck:} dynamic context construction is expensive
\item \textbf{Context Sensitivity:} performance depends on example selection
\item \textbf{LLM Bias:} encoder may overfit to current LLM behavior
\end{itemize}

\section{Design for Future LLM Compatibility}

To ensure compatibility with future LLMs, we propose:

\begin{itemize}
\item separating representation $h$ from interface $z$
\item enforcing intrinsic structure via auxiliary losses
\item using structured prompts for stability
\item optionally training with multiple LLMs
\end{itemize}

The goal is to learn \textbf{interface-stable representations} that generalize across models.

\section{Implementation}

\subsection{Repository Structure}

The current codebase separates preprocessing, training-time loading, and analysis utilities:
\begin{itemize}
\item \texttt{code/data\_cli.py}: inspect RelBench datasets through DuckDB
\item \texttt{code/preprocess\_train.py}: materialize train data and save reusable history artifacts
\item \texttt{code/train\_data.py}: core train-data utilities, history construction, and cache loading
\item \texttt{code/train.py}: thin runtime entrypoint for iterating over preprocessed batches
\item \texttt{code/configs/train\_config.py}: experiment configuration
\end{itemize}

\textbf{Pros:} clear separation of responsibilities; expensive preprocessing can be reused across runs; training code stays small.\\
\textbf{Cons:} more moving parts; preprocessing and runtime must remain schema-compatible.

\subsection{Dataset and Task Loading}

The implementation uses RelBench as the source of dataset and task definitions. Raw relational tables are materialized through the RelBench cache, while tasks are loaded through the task registry. DuckDB is used as a practical relational access layer for inspection and for compatibility with future non-RelBench datasets.

\textbf{Pros:} reuses well-defined benchmark tasks; keeps access to the raw relational database; compatible with future extensions.\\
\textbf{Cons:} some datasets/tasks are expensive to load; the implementation inherits RelBench task assumptions and timestamp schedules.

\subsection{Preprocessing-First Training Pipeline}

Rather than recomputing train histories inside every training iteration, the repository now uses an explicit preprocessing stage:
\begin{enumerate}
\item load each configured dataset and task
\item materialize the task train table
\item compute valid prior history candidates for each train datapoint
\item save train tables, history candidate caches, and a manifest to disk
\end{enumerate}

At runtime, \texttt{train.py} only loads these preprocessed artifacts and constructs batches.

\textbf{Pros:} much lower runtime overhead; reusable across epochs and runs; easier debugging of data issues.\\
\textbf{Cons:} extra preprocessing step; cached data can become stale if task semantics or config change.

\subsection{History Source Options}

The code supports two history sources:
\begin{itemize}
\item \texttt{task\_table}: sample history from earlier rows already present in the task train table
\item \texttt{dataset}: build history from task-valid prior timestamps derived from the raw dataset and task definition
\end{itemize}

\paragraph{Task-table history.}
This mode samples earlier labeled rows directly from the train split of the task table.

\textbf{Pros:} simple; fast; every sampled history item is already a valid labeled task example.\\
\textbf{Cons:} sparse; limited to timestamps already materialized by the task; less suitable for large-context ICL.

\paragraph{Dataset-based history.}
This mode uses the task's natural train-time schedule, computes valid labels from the raw database, and then samples prior valid history candidates.

\textbf{Pros:} more faithful to the intended relational/temporal setting; better suited for large history length $k$; more compatible with future ICL-style context construction.\\
\textbf{Cons:} more expensive to preprocess; task validity constraints may make candidate construction nontrivial.

\subsection{Caching Options}

Two cache layers are implemented:
\begin{itemize}
\item \texttt{cache\_dataset\_history\_labels}: in-memory timestamp-level label cache for dataset-based history
\item \texttt{cache\_train\_history\_candidates}: saved candidate-history pools for each train datapoint
\end{itemize}

The second cache is the main reusable cache for training. It stores compact history candidate references, not fully materialized sampled histories.

\textbf{Pros:} reduces repeated task recomputation; reusable across batches, epochs, and later runs; compact relative to saving full histories.\\
\textbf{Cons:} can still be large for dense tasks; requires cache invalidation discipline when configuration changes.

\subsection{Parallel History Construction}

At runtime, history construction can use one of two CPU strategies:
\begin{itemize}
\item \texttt{grouped\_vectorized}: group batch items by $(\text{dataset}, \text{task})$ and process them in a single process with lower Python overhead
\item \texttt{multiprocess}: distribute grouped history construction across CPU worker processes
\end{itemize}

\textbf{Grouped/vectorized mode.}
\textbf{Pros:} low overhead; good cache locality; often fastest for moderate batch sizes.\\
\textbf{Cons:} still limited by single-process Python throughput.

\textbf{Multiprocess mode.}
\textbf{Pros:} exploits independent per-item history construction; may help for very large batches.\\
\textbf{Cons:} process-management overhead; benefit depends on batch composition and memory-sharing behavior.

\subsection{Current Limitations of the Implementation}

The current repository stops at data preparation and mixed-batch construction. It does not yet implement:
\begin{itemize}
\item temporal graph encoders
\item embedding generation for history items
\item the LLM projection/interface module
\item end-to-end optimization against the frozen LLM
\end{itemize}

Thus, the implementation should currently be viewed as a \textbf{data and training-example construction layer} for the larger RFM pipeline, not the full model stack.

\subsection{Current Zero-Shot Inference Implementation}

The repository now also contains a zero-shot inference path under \texttt{code/zero\_shot.py}. Unlike the preprocessing-first training pipeline, this path constructs history and graph context online at inference time and directly queries a frozen local LLM.

\paragraph{Main supporting files.}
\begin{itemize}
\item \texttt{code/zero\_shot.py}: runtime entrypoint, CLI, streaming, logging, and MAE computation
\item \texttt{code/zero\_shot\_llm.py}: local LLM wrapper, model-path registry, and batched generation helper
\item \texttt{code/inference\_history.py}: inference-time history construction, validation, and history samplers
\item \texttt{code/grag.py}: GraphRAG store, temporal retrieval, optional semantic retrieval, and prompt builder
\item \texttt{code/configs/inference\_config.py}: single-dataset, single-task inference configuration
\end{itemize}

\paragraph{High-level runtime.}
For each evaluation query, the implementation now performs:
\begin{enumerate}
\item load the dataset and task through RelBench
\item materialize a raw relational database for inference
\item build a temporal GraphRAG index over the database
\item construct task-valid prior history for the query entity
\item retrieve $k$-hop graph neighborhoods under temporal cutoff constraints
\item assemble a prompt for a local LLaMA model
\item decode a scalar prediction
\item compute MAE when labels are available
\end{enumerate}

\subsection{Inference Query Semantics}

Each inference item is a task row with:
\begin{itemize}
\item entity identifier
\item query timestamp
\item optional target label
\end{itemize}

For a query:
\begin{equation}
(x,t) \rightarrow y
\end{equation}

the runtime first attempts to build a prior context:
\begin{equation}
\mathcal{C}_{hist} = \{(x,t_i,y_i)\}_{i=1}^{k}, \quad t_i < t
\end{equation}

When the history source is dataset-based, a stronger temporal validity rule is enforced:
\begin{equation}
t_i + \Delta_{task} \le t
\end{equation}

where $\Delta_{task}$ is the task horizon used to define the label window in RelBench. This prevents a history item's label from overlapping the final query's future interval.

\subsection{Inference-Time History Construction}

The current zero-shot implementation supports two sources of history:
\begin{itemize}
\item \texttt{task\_table}
\item \texttt{dataset}
\end{itemize}

\paragraph{Task-table history.}
Earlier labeled rows already present in the task table are sampled directly.

\paragraph{Dataset-based history.}
This is the more general current mode. For each query:
\begin{enumerate}
\item candidate timestamps are collected from raw tables with valid time columns
\item only timestamps strictly before the query time are retained
\item \texttt{task.make\_table(db, candidate\_timestamps)} is used to materialize task-valid rows
\item the resulting rows are filtered to the same entity
\item rows violating the non-overlap horizon constraint are removed
\item duplicate rows on entity and timestamp are removed
\item a history sampler chooses up to \texttt{history\_length} rows
\end{enumerate}

This design makes zero-shot history faithful to the task semantics rather than merely collecting arbitrary raw rows involving the entity.

\subsection{History Sampling Strategies}

Three history sampling strategies are currently implemented:
\begin{itemize}
\item \texttt{most\_recent\_k}
\item \texttt{random\_prior}
\item \texttt{recent\_min\_overlap}
\end{itemize}

\paragraph{Most recent $k$.}
Select the most recent valid prior rows.

\paragraph{Random prior.}
Uniformly sample from valid prior rows.

\paragraph{Recent min-overlap.}
Prefer recent valid rows while trying to keep selected timestamps separated by at least one task horizon before filling remaining slots.

The current default in the code is deterministic and does not require labels for tuning:
\begin{equation}
\texttt{history\_sampling\_strategy} = \texttt{most\_recent\_k}
\end{equation}

\subsection{GraphRAG Structure Used at Inference}

The GraphRAG store built in \texttt{code/grag.py} contains:
\begin{itemize}
\item one node for each primary-key row in each table
\item typed edges for FK$\rightarrow$PK relations
\item reverse adjacency for backward traversal
\item a task-entity history index
\item static row references for non-temporal rows
\end{itemize}

Thus retrieval is not limited to one direction of the schema graph. It can traverse both explicit FK edges and their reverse counterparts while still enforcing temporal cutoffs.

\subsection{Context Retrieved for Each Prompt Entry}

For each prompt entrypoint, the current implementation constructs:
\begin{itemize}
\item static rows for the entity
\item prior history rows before the entrypoint cutoff
\item $k$-hop neighbors of history rows
\item $k$-hop neighbors of the query entity node itself
\item optional same-table similar-entity evidence when semantic retrieval is enabled
\end{itemize}

An important current design choice is that the query entity's own neighborhood is always considered. Therefore, even if the task-history list is empty, the runtime still attempts to retrieve the query node's own temporal relational neighborhood before the cutoff.

\subsection{Temporal Safety of Retrieved Context}

The present implementation attempts to maintain zero-shot temporal validity through the following rules:
\begin{itemize}
\item history rows must occur strictly before the query time
\item dataset-derived history rows must satisfy the non-overlap horizon condition
\item neighbor retrieval is cutoff-aware
\item same-timestamp neighborhood edges may be included when they belong to the already-allowed row snapshot
\end{itemize}

This is intended to preserve the contemporaneous relational neighborhood of valid history rows while excluding truly future information.

\subsection{Prompt Structure}

The prompt assembled in \texttt{code/grag.py} currently contains:
\begin{enumerate}
\item task header
\item strict instruction to return only a number
\item one short explanation of what kinds of context are included
\item static information block
\item zero or more prior history examples with labels
\item history and graph-neighborhood context for those examples
\item the final unlabeled query block
\item short direct prediction instruction
\item \texttt{Answer:}
\end{enumerate}

The design goal is to simulate in-context learning for the same task and entity: prior history examples act as demonstrations, while the final query asks the LLM to infer the missing target.

\subsection{Neighbor Rendering}

All retrieved neighbors are currently included in the prompt. The implementation does not do prompt-side truncation beyond the retrieval limits already imposed by:
\begin{itemize}
\item \texttt{num\_hops}
\item \texttt{top\_k}
\end{itemize}

Neighbors are rendered in a graph-aware style that tries to make the relation structure visible to the LLM rather than presenting them as an unstructured list of rows.

\subsection{Optional Semantic Retrieval}

The zero-shot CLI also exposes an optional retrieval mode beyond explicit PK/FK structure. When enabled, the system retrieves same-table semantically similar entities using generic static-attribute similarity. For each similar entity, the prompt may include:
\begin{itemize}
\item entity identity
\item heuristic similarity score
\item static summary
\item pre-cutoff history rows
\end{itemize}

This path is disabled by default and is intended as a generic cold-start fallback rather than a task-specific learned module.

\subsection{Local LLM Interface}

The current inference runtime uses locally downloaded LLaMA checkpoints through Hugging Face \texttt{transformers}. The wrapper currently supports:
\begin{itemize}
\item \texttt{1b}
\item \texttt{3b}
\item \texttt{8b}
\end{itemize}

Generation is decoder-only, deterministic, and followed by numeric extraction from the returned text.

\subsection{Two-Stage Runtime and Streaming}

The runtime is structured as a two-stage system:
\begin{enumerate}
\item CPU-side context generation
\item GPU-side LLM generation
\end{enumerate}

Prepared prompts are streamed to the LLM as soon as they become available. Thus the LLM does not need to wait for the entire context batch to finish if:
\begin{itemize}
\item prompt preparation is producing ready items
\item the LLM batch size is small
\end{itemize}

\subsection{Batching and Parallelism}

The current zero-shot implementation exposes:
\begin{itemize}
\item context-batch size
\item LLM-batch size
\item worker count for example-level prompt preparation
\item worker count for frontier expansion inside multi-hop graph search
\end{itemize}

Parallelism currently exists at multiple levels:
\begin{itemize}
\item across examples in a context batch
\item within prompt construction across entrypoints
\item within a hop frontier during graph expansion
\end{itemize}

This may improve throughput, but excessive nested parallelism can also increase overhead and harm latency. Therefore, the current runtime should be viewed as an exploratory parallel implementation rather than a final optimized scheduler.

\subsection{Evaluation}

When labels are available, the runtime computes MAE:
\begin{equation}
\mathrm{MAE} = \frac{1}{N}\sum_{i=1}^{N} |\hat{y}_i - y_i|
\end{equation}

If labels are absent, predictions are still produced but MAE is skipped.

\subsection{Current Limitations of the Zero-Shot Path}

The present implementation remains expensive for several reasons:
\begin{itemize}
\item dataset-based history construction repeatedly materializes task-valid rows
\item multi-hop graph retrieval can expand rapidly
\item prompt size grows quickly when all neighbors are included
\item semantic retrieval remains heuristic
\item throughput and latency depend strongly on worker settings
\end{itemize}

Nevertheless, unlike the earlier training-only path, the current repository now includes a full end-to-end zero-shot inference stack:
\begin{itemize}
\item online history construction
\item temporal GraphRAG retrieval
\item prompt assembly
\item local LLM generation
\item metric computation
\end{itemize}

\section{Evaluation Plan}

\subsection{Baselines}
\begin{itemize}
\item TGNN + MLP
\item RDBLearn
\item TabPFN
\end{itemize}

\subsection{Metrics}
\begin{itemize}
\item predictive performance (MAE, AUC, MRR)
\item inference latency
\item sampling overhead
\end{itemize}

\subsection{Ablations}
\begin{itemize}
\item with/without history
\item random vs structured sampling
\item projection vs no projection
\item auxiliary loss vs none
\end{itemize}

\section{Conclusion}

We introduce a new paradigm for relational learning:

\begin{quote}
\textit{Train relational encoders not for direct prediction, but for compatibility with in-context reasoning in frozen LLMs.}
\end{quote}

This approach opens a path toward relational foundation models that are modular, generalizable, and future-proof.

\bibliographystyle{plain}
\bibliography{refs}

\end{document}
